\section{Formal verification of closed-system and semidefinite consequences}
\label{sec:formal-consequences}

The second Lean package verifies consequences of the certificate theorem
and two examples that delimit them. It uses the actual closed
set $T=\{x:f_i(x)\leq0\text{ for every }i\}$ and the ordinary convex hull.
Its Shor projection is
\[
 P=\left\{x:\exists X,\quad
 \begin{pmatrix}1&x\transpose\\x&X\end{pmatrix}\succeq0,
 \quad \tr(A_iX)+2b_i\transpose x+c_i\leq0\text{ for every }i\right\}.
\]
The formalization proves the Schur-complement equivalence with covariance
slack $Y\succeq0$, where $X=xx\transpose+Y$ and the lifted inequalities
are $f_i(x)+\tr(A_iY)\leq0$. It also proves $S\subseteq T\subseteq P$
and convexity of $P$.

\paragraph{Consequences and their assumptions.}
If $S\ne\varnothing$ and asymptotic hyperplane convexity holds, then
\[
 \conv T\ne\R^n
 \quad\Longleftrightarrow\quad
 \exists\lambda\geq0,\ \lambda\ne0:
 A_\lambda\succeq0,\ (A_\lambda,b_\lambda)\ne(0,0),
\]
and
\[
 P=\R^n\quad\Longleftrightarrow\quad
 \conv S=\R^n\quad\Longleftrightarrow\quad\conv T=\R^n.
\]
Here asymptotic hyperplane convexity means that, for every nonzero normal,
the homogeneous image is convex on arbitrarily distant sweeping hyperplanes;
HHC implies this condition. A PSD aggregate is nonpositive throughout $P$,
because the trace pairing of two PSD matrices is nonnegative. A nonconstant
PSD quadratic is unbounded above, so its nonpositive sublevel set is proper.
These facts prove the certificate-to-properness implications without any
feasibility or hyperplane-convexity premise.

\paragraph{An unconditional Shor characterization.}
With strict feasibility alone, the package proves
\[
 P=\R^n
 \quad\Longleftrightarrow\quad
 \text{every nonnegative PSD aggregation has }A_\lambda=b_\lambda=0.
\]
The verified converse uses open-set separation and does not require a
closed linear image of the PSD cone. For fixed $x$, suppose there were no
$Y\succeq0$ satisfying all $f_i(x)+\tr(A_iY)<0$. Separate the convex affine
image of PSD matrices from the open negative orthant. This gives nonzero
$\lambda\geq0$ with
\[
 0\leq f_\lambda(x)+\tr(A_\lambda Y)\qquad(Y\succeq0).
\]
Scaling the rank-one matrices $Y=vv\transpose$ proves $A_\lambda\succeq0$,
and setting $Y=0$ gives $f_\lambda(x)\geq0$. If every such aggregation is
trivial, evaluation at a strict feasible point instead gives
$f_\lambda(x)=c_\lambda<0$, a contradiction. Thus every fiber admits an
actual PSD slack with strict inequalities, which proves membership in the
ordinary Shor projection, not only in its closure.

\paragraph{Exact finite SDP characterization.}
Normalize multipliers by $\lambda\geq0$ and $\sum_i\lambda_i=1$, and
require $A_\lambda\succeq0$. The resulting set is compact or empty.
There are $n(n+1)/2+n$ independent coordinates of $(A_\lambda,b_\lambda)$.
The formalization proves that every signed coordinate objective attains
its maximum on a nonempty normalized set. All these coordinates vanish
throughout that set exactly when every one of the $n(n+1)+2n$ signed
programs has optimum zero, or the common feasible set is empty. The
normalization argument connects this to the full multiplier cone and
the preceding whole-space characterizations. These are mathematical
equivalences about exact feasibility and exact optimum values. They do
not verify floating-point SDP solvers or turn a near-zero objective into
a zero certificate.

\paragraph{Verified boundary examples.}
The first example uses
\[
 f_1=x_1x_2-1,\qquad f_2=1-x_1x_2,\qquad f_3=x_1^2-x_2^2
 \quad\text{on }\R^3.
\]
Its strict system is empty, its closed set is nonempty, and its closed
feasible hull lies in $|x_1|\leq1$. Every PSD aggregation has zero quadratic
and linear parts. Nevertheless HHC holds. Thus strict feasibility cannot
be removed from the closed-system consequence.

The second example uses $f_1=x_1^2-4$ and
$f_2=-x_1^2+2x_1+3$ on $\R^3$. Its strict feasible set is the convex strip
$-2<x_1<-1$, and HHC holds. Every nonzero nonnegative globally convex
aggregation is strictly negative at zero. Zero also belongs to the Shor
projection: $X=\operatorname{diag}(7/2,0,0)$ gives a PSD block and both
lifted residuals equal $-1/2$. Hence that projection is proper but strictly
larger than the hull. The examples include proofs of HHC: a formal algebraic
proof of Dines' two-form theorem establishes HHC for pairs of forms, and a
linear-image argument handles the repeated negative form.

\paragraph{Scope.}
The proof files extend \texttt{Formal.QuadraticAggregation}. The frozen
claims, declaration coverage, and audit inputs are in
\path{supplement/lean/audits/28-quadratic-aggregation-consequences/}.
The full good-aggregation hull theorem, the general sufficient criteria
based on classical three-form convexity, additional examples, numerical
solver correctness, and literature priority are outside this package.
The infinite-aggregation construction has a separate verification scope.
